\section{Physical excitation counts and replica energy}
\label{sec:replica_defect}

The physical gap and defect-mass calculation in
\cite[Section~7]{OpenAI2026AreaLaw} take place on the physical replicas,
with identity operators on the auxiliary factors.

\begin{definition}[Physical replica operators]
\label{def:replica_operators}
\lean{Matrix.replicaHamiltonian,Matrix.replicaDefectCount}
\leanok
Let $\mathcal H=\mathbb C^A$, let $H$ be an operator on $\mathcal H$,
and let $\Omega\in\mathcal H$. For $k\ge0$, define
\begin{align}
 H_\Sigma&=\sum_{i=0}^{k-1} I^{\otimes i}\otimes H\otimes I^{\otimes(k-1-i)},
 \label{eq:replica_hamiltonian}
\end{align}
\begin{align}
 N&=\sum_{i=0}^{k-1} I^{\otimes i}\otimes
       (I-|\Omega\rangle\langle\Omega|)\otimes I^{\otimes(k-1-i)}.
 \label{eq:replica_defect}
\end{align}
For $k>0$, write $\overline H=k^{-1}H_\Sigma$.
\end{definition}

\begin{theorem}[Product eigenvectors of the replica sum]
\label{thm:replica_product_eigenvector}
\lean{Matrix.replicaHamiltonian_mulVec_prod}
\leanok
\uses{def:replica_operators}
If $H\Omega=E_0\Omega$ with $E_0\in\mathbb C$, then
$H_\Sigma\Omega^{\otimes k}=kE_0\Omega^{\otimes k}$ for every $k\ge0$.
No normalization of $\Omega$ is required.
\end{theorem}
\begin{proof}
\leanok
Each summand in (\ref{eq:replica_hamiltonian}) sends the product vector to
$E_0\Omega^{\otimes k}$. Summing the $k$ summands proves the identity.
\end{proof}

\begin{theorem}[Physical mean energy with auxiliary factors]
\label{thm:replica_auxiliary_eigenvector}
\lean{Matrix.replicaMeanHamiltonian_kronecker_mulVec_prod}
\leanok
\uses{def:replica_operators}
Suppose $H\Omega=E_0\Omega$, where $E_0\in\mathbb R$, and $k>0$.
For every finite-dimensional auxiliary space $\mathcal K$ and every
$w\in\mathcal K$,
$(\overline H\otimes I_{\mathcal K})(\Omega^{\otimes k}\otimes w)
=E_0\Omega^{\otimes k}\otimes w$.
Neither vector is required to be normalized.
\end{theorem}
\begin{proof}
\leanok
\uses{thm:replica_product_eigenvector}
The identity on $\mathcal K$ fixes $w$; divide the product eigenvalue by $k$.
\end{proof}

\begin{theorem}[The one-copy gap on physical replicas]
\label{thm:replica_gap}
\lean{Matrix.PosSemidef.replica_gap}
\leanok
\uses{def:replica_operators}
Let $E_0,g\in\mathbb R$. If
$H-E_0I-g(I-|\Omega\rangle\langle\Omega|)$ is positive semidefinite, then
$H_\Sigma-kE_0I-gN$ is positive semidefinite for every $k\ge0$.
This assertion requires neither normalization nor an eigenvector equation.
\end{theorem}
\begin{proof}
\leanok
Tensor the one-copy positive semidefinite operator with identities on all
other copies and sum over its location. Linearity gives
$H_\Sigma-kE_0I-gN$.
\end{proof}

\begin{theorem}[Positivity of the physical defect count]
\label{thm:replica_defect_positive}
\lean{Matrix.posSemidef_replicaDefectCount}
\leanok
\uses{def:replica_operators}
If $\|\Omega\|=1$, then $N\ge0$ for every $k\ge0$.
\end{theorem}
\begin{proof}
\leanok
The Cauchy--Schwarz inequality gives
$|\Omega\rangle\langle\Omega|\le I$. Tensoring each complement with
identities and summing preserves positivity.
\end{proof}

\begin{theorem}[Actual cutoff mass from the physical mean energy]
\label{thm:replica_cutoff_mass}
\lean{Matrix.spectralCutoff_replica_gap_mass_ge}
\leanok
\uses{def:replica_operators}
Assume $\|\Omega\|=1$, $g>0$, $k>0$, $\tau>0$, and
$H-E_0I-g(I-|\Omega\rangle\langle\Omega|)\ge0$, with $E_0\in\mathbb R$.
Let $v\in\mathcal H^{\otimes k}\otimes\mathcal K$ be a unit vector, and
let $\Pi=\mathbf1_{(-\infty,\tau k]}(N)$ be the actual spectral projection.
Then
\begin{align}
 \langle v,(\Pi\otimes I_{\mathcal K})v\rangle
 &\ge 1-\frac{\langle v,(\overline H\otimes I_{\mathcal K})v\rangle-E_0}
                  {g\tau}.
 \label{eq:replica_cutoff_mass}
\end{align}
The inner products in (\ref{eq:replica_cutoff_mass}) are real because the
operators are self-adjoint.
\end{theorem}
\begin{proof}
\leanok
\uses{thm:replica_gap,thm:replica_defect_positive,thm:spectral_cutoff_mass}
The spectral first-moment inequality gives
$\langle v,(\Pi\otimes I)v\rangle\ge
1-\langle v,(N\otimes I)v\rangle/(\tau k)$.
The replica gap gives
$g\langle v,(N\otimes I)v\rangle
\le\langle v,(H_\Sigma\otimes I)v\rangle-kE_0$.
Substitution proves (\ref{eq:replica_cutoff_mass}).
\end{proof}
